\begin{frame}<1-3>[t]{新知探索}
\begin{campuscanvas}
% Source slide 15, objects 25; visible 2-3
\only<2-3|handout:1>{\node[anchor=north west,inner sep=0,text=black] at (70.000,150.000) {\begin{minipage}[t]{142.500mm}\raggedright\fontsize{16}{24.00}\selectfont \textbf{6. 表示集合的方法：描述法}\end{minipage}};}
% Source slide 15, objects 24; visible 3
\only<3|handout:1>{\node[anchor=north west,inner sep=0,text=black] at (70.000,250.000) {\begin{minipage}[t]{142.500mm}\raggedright\fontsize{13}{19.50}\selectfont 无限集一般不能用列举法表示。有限集如果元素太多或叫不出名字来，例如某池塘里所有鱼的集合，也不便用列举法来表示。\par 这时可以\red{把集合中元素共有的，也只有该集合中元素才有的属性描述出来，以确定这个集合}。这种表示法叫作\red{描述法}。\par 集会时介绍嘉宾常用列举法，致辞里说“女士们，先生们”用的就是描述法。\end{minipage}};}
\end{campuscanvas}
\end{frame}
